% ============================================================
%  CHAPTER 4: IMPLEMENTATION DETAILS
% ============================================================

\chapter{IMPLEMENTATION DETAILS}
\label{chap:implementation}
\label{chap:deployment}

\section{System Architecture and Repository Layout}
\label{sec:repo_layout_sec}

The QuantumSafe VPN codebase is organized into two primary sub-projects: the native Kotlin Android mobile client and the companion Python Linux gateway server (Table~\ref{tab:repo_layout}).

\begin{table}[htbp]
  \centering
  \scriptsize
  \setlength{\tabcolsep}{3pt}
  \caption{Project Codebase Layout and Key Component Responsibilities.}
  \label{tab:repo_layout}
  \begin{tabularx}{\linewidth}{@{}p{2.6cm}p{3.0cm}Y@{}}
    \toprule
    \textbf{Subdirectory / Path} & \textbf{Component Name} & \textbf{Primary Engineering Responsibility} \\
    \midrule
    \texttt{android/app/}\allowbreak\texttt{src/.../vpn} & \texttt{QuantumSafe}\allowbreak\texttt{VpnService.kt} & Manages TUN interface lifecycle, socket protection, and packet loops. \\
    \texttt{android/app/.../}\allowbreak\texttt{crypto} & \texttt{PqcEngine.kt} & Pure-Kotlin ML-KEM, ML-DSA, Keccak SHA3-256, and AES-GCM engine. \\
    \texttt{android/app/.../ui} & \texttt{HomeScreen, Theme.kt} & Jetpack Compose declarative UI screens and minimal white design tokens. \\
    \texttt{android/app/.../}\allowbreak\texttt{storage} & \texttt{Security}\allowbreak\texttt{Manager.kt} & Hardware-backed Android Keystore configuration encryption. \\
    \texttt{vpn/tunnel\_}\allowbreak\texttt{server.py} & Linux Gateway Daemon & Multi-client TCP listener, PQC handshake dispatch, and TUN routing. \\
    \texttt{vpn/ip\_pool.py} & Dynamic IP Lease Manager & Thread-safe allocation and reclamation of 10.8.0.0/24 client virtual IPs. \\
    \texttt{pqc/mlkem.py,}\allowbreak\texttt{mldsa.py} & Python PQC Modules & ML-KEM-1024 encapsulation/decapsulation and ML-DSA-65 identity signing. \\
    \texttt{crypto/}\allowbreak\texttt{classical.py} & Classical Cryptography & HKDF-SHA256 session key expansion and AES-256-GCM AEAD framing. \\
    \bottomrule
  \end{tabularx}
\end{table}

\section{Android VpnService Lifecycle and Socket Protection}
\label{sec:android_vpn_lifecycle}

The Android client extends \texttt{android.net.VpnService}. Under Android 14 (API Level 34), foreground services must declare explicit types in the app's \texttt{AndroidManifest.xml}:
\begin{lstlisting}[language=XML, basicstyle=\ttfamily\scriptsize]
<service android:name=".QuantumSafeVpnService"
    android:permission="android.permission.BIND_VPN_SERVICE"
    android:foregroundServiceType="systemExempted"
    android:exported="false">
    <intent-filter><action android:name="android.net.VpnService" /></intent-filter>
</service>
\end{lstlisting}
The \texttt{systemExempted} foreground type prevents OS battery optimization daemons from freezing the tunnel transport during long periods of background execution.

\subsection{Kernel-Level Socket Protection Bypass}
To prevent an infinite routing loop where encrypted tunnel packets are routed back into the virtual TUN interface, \texttt{QuantumSafeVpnService} invokes the native \texttt{protect()} system call prior to initiating TCP transport:
\begin{lstlisting}[language=Kotlin, basicstyle=\ttfamily\scriptsize]
val socket = Socket()
if (!protect(socket)) {
    throw IllegalStateException("Critical failure: Failed to protect transport socket!")
}
socket.tcpNoDelay = true // Disable Nagle's algorithm for low-latency transmission
socket.connect(InetSocketAddress(serverHost, serverPort), 10000)
\end{lstlisting}
Enabling \texttt{tcpNoDelay} eliminates artificial 40\,ms buffering delays for interactive frames.

\section{Pure-Kotlin Post-Quantum Handshake Engine}
\label{sec:kotlin_pqc_handshake}

Listing~\ref{lst:handshake_pqc_engine} illustrates the execution of the 4-stage post-quantum handshake inside \texttt{PqcEngine.kt} over the protected transport socket.

\begin{lstlisting}[language=Kotlin, caption={4-Stage Post-Quantum Handshake Execution in PqcEngine.kt.}, label={lst:handshake_pqc_engine}, basicstyle=\ttfamily\scriptsize]
fun performHandshake(socket: Socket): HandshakeResult {
    val inputStream = DataInputStream(socket.getInputStream())
    val outputStream = DataOutputStream(socket.getOutputStream())

    // Stage 1: Receive ML-DSA-65 Server Public Key (1952 Bytes)
    val serverDsaPk = readFrame(inputStream)
    if (serverDsaPk.size != 1952) throw SecurityException("Invalid ML-DSA key size")

    // Stage 2: Generate Ephemeral ML-KEM-1024 Keypair & Transmit Public Key (1568 Bytes)
    val (kemPk, kemSk) = generateMlkemKeypair()
    sendFrame(outputStream, kemPk)

    // Stage 3: Receive Ciphertext (1568 Bytes) and Server Signature (3309 Bytes)
    val ciphertext = readFrame(inputStream)
    val signature = readFrame(inputStream)
    val transcript = "QuantumSafeVPN-v1".toByteArray(Charsets.UTF_8) + kemPk + ciphertext
    if (!verifyMldsa(serverDsaPk, transcript, signature)) {
        throw SecurityException("Server identity verification failed! Invalid signature.")
    }

    // Decapsulate Shared Secret & Receive Assigned IP Configuration
    val sharedSecret = decapsulateMlkem(kemSk, ciphertext)
    val config = parseConfigJson(readFrame(inputStream))
    val sessionKey = deriveSessionKey(sharedSecret) // HKDF-SHA256 expansion

    return HandshakeResult(sessionKey, config)
}
\end{lstlisting}

\section{Pure-Kotlin Keccak-$f[1600]$ Sponge Implementation}
\label{sec:keccak_impl_sec}

To guarantee bit-level parity with Python's \texttt{hashlib.sha3\_256}, \texttt{PqcEngine} incorporates a pure-Kotlin implementation of the NIST FIPS 202 Keccak sponge permutation:
\begin{lstlisting}[language=Kotlin, basicstyle=\ttfamily\scriptsize]
private fun keccakF1600(state: LongArray) {
    for (round in 0 until 24) {
        val c = LongArray(5) { x ->
            state[x] xor state[x+5] xor state[x+10] xor state[x+15] xor state[x+20]
        }
        for (x in 0 until 5) {
            val d = c[(x + 4) % 5] xor java.lang.Long.rotateLeft(c[(x + 1) % 5], 1)
            for (y in 0 until 5) state[x + 5 * y] = state[x + 5 * y] xor d
        }
        val b = LongArray(25) { i ->
            val x = i % 5; val y = i / 5
            java.lang.Long.rotateLeft(state[x + 5 * y], KECCAK_ROTATION_CONSTANTS[x][y])
        }
        for (x in 0 until 5) {
            for (y in 0 until 5) state[x + 5 * y] = b[x + 5 * y] xor (b[(x + 1) % 5 + 5 * y].inv() and b[(x + 2) % 5 + 5 * y])
        }
        state[0] = state[0] xor KECCAK_ROUND_CONSTANTS[round]
    }
}
\end{lstlisting}

\section{Bidirectional Forwarding and Packet Framing}
\label{sec:bidirectional_forwarding}

Once \texttt{VpnService.Builder} establishes the virtual \texttt{tun0} interface, two concurrent coroutine workers are launched on \texttt{Dispatchers.IO}. Listing~\ref{lst:uplink_worker} illustrates the uplink forwarding pipeline.

\begin{lstlisting}[language=Kotlin, caption={Uplink Packet Forwarding and AEAD Encapsulation Worker.}, label={lst:uplink_worker}, basicstyle=\ttfamily\scriptsize]
private fun CoroutineScope.launchUplinkWorker(pfd: ParcelFileDescriptor, socket: Socket, pqcEngine: PqcEngine) = launch(Dispatchers.IO) {
    val inputStream = FileInputStream(pfd.fileDescriptor)
    val outputStream = DataOutputStream(socket.getOutputStream())
    val buffer = ByteArray(32767)

    while (isActive && isRunning) {
        val bytesRead = inputStream.read(buffer)
        if (bytesRead > 0) {
            val rawPacket = buffer.copyOf(bytesRead)
            val seq = sendSequenceCounter.incrementAndGet()
            val encryptedFrame = pqcEngine.encryptPacket(rawPacket, seq) // AES-GCM + IV + Tag

            synchronized(outputStream) {
                outputStream.writeInt(encryptedFrame.size)
                outputStream.write(encryptedFrame)
                outputStream.flush()
            }
            totalBytesSent.addAndGet(encryptedFrame.size.toLong())
        }
    }
}
\end{lstlisting}

\section{Python Linux Gateway Daemon Implementation}
\label{sec:gateway_daemon_impl}

On the server, \texttt{vpn/tunnel\_server.py} manages incoming client connections using Python's \texttt{threading} architecture. Listing~\ref{lst:python_server_handler} details client dispatching and sequence verification.

\begin{lstlisting}[language=Python, caption={Python Linux Gateway Client Handler and Replay Validation.}, label={lst:python_server_handler}, basicstyle=\ttfamily\scriptsize]
def handle_client_connection(tun, conn, addr):
    try:
        session_key, assigned_ip = perform_handshake(conn, addr)
        session_router.register_client(assigned_ip, conn, session_key)
        last_seq = 0
        while True:
            encrypted_frame = recv_frame(conn)
            plaintext = decrypt(session_key, encrypted_frame) # AES-256-GCM verification
            seq = struct.unpack("!Q", plaintext[:8])[0]
            if seq <= last_seq:
                continue # Discard out-of-order or replayed packet
            last_seq = seq
            tun.write_packet(plaintext[8:]) # Forward IPv4 packet into Linux tun0
    finally:
        session_router.unregister_client(assigned_ip)
        ip_pool.release(assigned_ip)
        conn.close()
\end{lstlisting}

The kernel routes outbound packets via \texttt{iptables -t nat -A POSTROUTING -s 10.8.0.0/24 -o eth0 -j MASQUERADE}. Production APKs are compiled using Gradle 8.9 and R8 shrinking, yielding an optimized 13.4\,MB binary.

\section{Production Deployment Architecture and Gateway Provisioning}
\label{sec:production_deployment_provisioning}

Deploying QuantumSafe VPN involves establishing a hardened infrastructure consisting of the Linux Gateway and Android Endpoints:
\begin{lstlisting}[language=bash, basicstyle=\ttfamily\scriptsize]
# Enable runtime and persistent kernel IPv4 forwarding
sudo sysctl -w net.ipv4.ip_forward=1
echo "net.ipv4.ip_forward = 1" | sudo tee -a /etc/sysctl.d/99-vpn-forward.conf

# Masquerade outbound traffic from VPN clients (10.8.0.0/24) via physical WAN (eth0)
sudo iptables -t nat -A POSTROUTING -s 10.8.0.0/24 -o eth0 -j MASQUERADE
sudo iptables -A FORWARD -i tun0 -o eth0 -j ACCEPT
sudo iptables -A FORWARD -i eth0 -o tun0 -m state --state RELATED,ESTABLISHED -j ACCEPT
sudo iptables -A INPUT -p tcp --dport 7000 -j ACCEPT
\end{lstlisting}

The systemd service daemon configuration at \texttt{/etc/systemd/system/quantumsafe-vpn.service} ensures resilient background persistence:
\begin{lstlisting}[language=bash, basicstyle=\ttfamily\scriptsize]
[Unit]
Description=QuantumSafe Post-Quantum VPN Server Daemon
After=network.target network-online.target

[Service]
Type=simple
User=root
WorkingDirectory=/opt/quantumsafe-vpn
Environment="PYTHONPATH=." "VPN_SERVER_PORT=7000"
ExecStartPre=/bin/bash /opt/quantumsafe-vpn/scripts/setup_network.sh
ExecStart=/usr/bin/python3 /opt/quantumsafe-vpn/vpn/tunnel_server.py
Restart=always
RestartSec=5s

[Install]
WantedBy=multi-user.target
\end{lstlisting}

\section{Android Client Deployment and Operational Troubleshooting}
\label{sec:troubleshooting_runbook}

The mobile client is installed on Android 8.0--14 devices via ADB sideloading or direct package installation:
\begin{lstlisting}[language=bash, basicstyle=\ttfamily\scriptsize]
adb install -r dist/QuantumSafe-VPN-release.apk
\end{lstlisting}
Upon first launch, the user configures the target gateway IP (\texttt{11.12.5.8}) and port (\texttt{7000}) in the Server Config tab, returns to Home, and taps Connect, accepting the system-level \texttt{VpnService} permission prompt.

Table~\ref{tab:troubleshooting_guide} details common operational issues, diagnostic indicators, and verified resolution procedures.

\begin{table}[htbp]
  \centering
  \scriptsize
  \setlength{\tabcolsep}{3pt}
  \caption{Operational Troubleshooting and Diagnostic Runbook.}
  \label{tab:troubleshooting_guide}
  \begin{tabularx}{\linewidth}{@{}p{3.2cm}p{3.8cm}Y@{}}
    \toprule
    \textbf{Observed Symptom} & \textbf{Probable Root Cause} & \textbf{Prescribed Resolution Procedure} \\
    \midrule
    Connection Timeout (UI ERROR) & Port 7000 blocked by firewall; daemon inactive & 1. Verify daemon: \texttt{ss -tulpn | grep 7000}.\\ & & 2. Allow TCP 7000 in firewall.\\ & & 3. Confirm network route to gateway IP. \\
    Connected; No Internet Access & Kernel IP forwarding off; NAT rule missing & 1. Check: \texttt{sysctl net.ipv4.ip\_forward}.\\ & & 2. Verify NAT: \texttt{iptables -t nat -L -v -n}.\\ & & 3. Run \texttt{sudo bash scripts/setup\_network.sh}. \\
    Signature Failure & Server key altered; MitM active & 1. Inspect server log for key renewal.\\ & & 2. Clear cached key in app Settings.\\ & & 3. Re-verify \texttt{server\_mldsa65\_public.key}. \\
    UI Freezes upon Connect & Socket loop bypass failed & 1. Confirm \texttt{protect(socket)} returned true.\\ & & 2. Ensure OS battery manager allows app.\\ & & 3. Restart application process. \\
    \bottomrule
  \end{tabularx}
\end{table}

\subsection{System Limitations and Maintenance Lifecycle}
\begin{enumerate}
    \item \textbf{TCP Encapsulation Overhead:} Tunneling IP packets over TCP can induce throughput degradation (TCP-over-TCP meltdown) on cellular links experiencing $> 5\%$ packet loss. A UDP/DTLS transport engine with automated TCP fallback is planned for subsequent iterations.
    \item \textbf{Single-Hop Star Topology:} The system connects clients to a centralized gateway; decentralized peer-to-peer mesh routing is not currently supported.
    \item \textbf{IPv4 Exclusivity:} The virtual TUN interface captures IPv4 traffic; native IPv6 dual-stack encapsulation remains an ongoing development objective.
    \item \textbf{Pure-Kotlin Mathematical Primitives:} While cross-platform and memory-safe, pure-Kotlin polynomial arithmetic does not leverage platform SIMD vector registers (ARM NEON), resulting in slightly higher CPU cycle counts than optimized C/assembly.
\end{enumerate}
The server enforces zero-logging (discarding packet payloads immediately and keeping session tables in volatile RAM), and server ML-DSA identity keys are rotated annually.
